%==============================================================================
%  CHAPTER 4 — IMPLEMENTATION AND DEPLOYMENT
%==============================================================================
\chapter{Implementation and Deployment}
\label{chap:implementation}

\section{Introduction}
This chapter details the realization of the design of
Chapter~\ref{chap:design}, sprint by sprint, following the Scrum plan of
Chapter~\ref{chap:context}. Each sprint section presents its backlog, the key
implementation points, and the delivered interface illustrated with application
screenshots. The chapter closes with the continuous-deployment chain and the
execution environments. The validation of all delivered increments is consolidated in
Chapter~\ref{chap:validation}.

% ===========================================================================
\section{Sprint 1 --- Foundation: Identity and Projects}
\label{sec:sprint1}
The first sprint lays the groundwork: a secure authentication system and the
project entity around which all later work is organized
(Figure~\ref{fig:sprint1-strip}).

\begin{figure}[H]\centering
  \includegraphics[width=\textwidth,keepaspectratio]{sprint1_features}
  \caption{Sprint 1 at a glance: authentication, projects, isolation}
  \label{fig:sprint1-strip}
\end{figure}

\begin{table}[H]\centering
\renewcommand{\arraystretch}{1.15}
\caption{Sprint~1 backlog}\label{tab:backlog-s1}
\small
\begin{tabularx}{\textwidth}{|c|X|c|}
\hline
\rowcolor{primary!15}
\textbf{\#} & \textbf{Task} & \textbf{Pts}\\
\hline
1 & Design the user and project data model & 2\\ \hline
1 & Implement signup / login / refresh with JWT & 3\\ \hline
1 & Build the login and signup user interface & 1\\ \hline
2 & Implement project CRUD with ownership checks & 2\\ \hline
2 & Bind each project to its own tracking experiment & 1\\ \hline
2 & Build the projects list and detail interface & 1\\ \hline
\end{tabularx}
\end{table}

\subsection{Authentication with JWT}
Authentication is stateless and based on \gls{jwt}. On login, the backend
verifies the password against its stored bcrypt hash and issues a short-lived
\emph{access} token and a longer-lived \emph{refresh} token. Protected
endpoints depend on a function that decodes the token and resolves the current
user; the refresh endpoint renews the session without re-entering credentials
(Figure~\ref{fig:screen_login}).

\bigshot{screen_login}{The authentication screen}{A returning user signs in
with email and password; new users switch to the sign-up form, and sessions
are kept alive transparently through the refresh token.}

\subsection{Project Management and Isolation}
A project is the organizing unit of the platform. Creating one also provisions
a dedicated experiment in the tracking server, so every run under the project
is grouped automatically. All project endpoints enforce an ownership check
that answers \texttt{404}, not \texttt{403}, for foreign resources: another
user's project simply does not exist for the caller. The frontend mirrors this
protection with an Angular route guard that redirects any unauthenticated
access to the login page (Figure~\ref{fig:screen_projects}).

\bigshot{screen_projects}{The projects workspace}{Each card is an isolated
project bound to its own tracking experiment, and the dialog creates a new
project; a user only ever sees the projects they own.}

\subsection{Sprint Review}
The increment was demonstrated to the supervisors and accepted: accounts,
sessions, and isolated project workspaces worked end to end on the cluster.
The retrospective confirmed the choice of stateless \gls{jwt} sessions and led
to one adjustment: answering \texttt{404} instead of \texttt{403} on foreign
resources, so that ownership is never leaked.

% ===========================================================================
\section{Sprint 2 --- Training and Tracking: the Engine}
\label{sec:sprint2}
The second sprint delivers the heart of the platform: uploading code and data,
running it as an isolated job on the cluster, and capturing the results
automatically (Figure~\ref{fig:sprint2-strip}).

\begin{figure}[H]\centering
  \includegraphics[width=\textwidth,keepaspectratio]{sprint2_features}
  \caption{Sprint 2 at a glance: upload, conversion, cluster runs, live logs}
  \label{fig:sprint2-strip}
\end{figure}

\begin{table}[H]\centering
\renewcommand{\arraystretch}{1.15}
\caption{Sprint~2 backlog}\label{tab:backlog-s2}
\small
\begin{tabularx}{\textwidth}{|c|X|c|}
\hline
\rowcolor{primary!15}
\textbf{\#} & \textbf{Task} & \textbf{Pts}\\
\hline
3 & Upload code/dataset to object storage & 3\\ \hline
3 & Auto-detect the entry script in an upload & 2\\ \hline
4 & Build the training pipeline (download $\to$ run $\to$ autolog) & 5\\ \hline
4 & Convert notebooks to scripts and inject the autolog wrapper & 3\\ \hline
5 & Enable automatic capture of metrics, parameters, artifacts & 5\\ \hline
6 & Stream live pipeline logs to the UI & 3\\ \hline
7 & Build the experiments browsing interface & 3\\ \hline
7 & Build the run-details view with metric cards and an artifact browser & 3\\ \hline
\end{tabularx}
\end{table}

\subsection{Upload and Entry-Script Detection}
Uploaded files are stored in the object store under a key derived from the
project and a short unique prefix. The backend supports scripts, notebooks,
and archives, and auto-detects the entry script by a fixed priority: an
explicit user choice, then conventional names (\texttt{main.py},
\texttt{train.py}), then the only script present, and finally the notebook
(Figure~\ref{fig:screen_upload}).

\bigshot{screen_upload}{The upload area}{Code (\texttt{.py} / \texttt{.ipynb} /
\texttt{.zip}) and datasets are dropped straight in the browser; the entry
script is detected automatically and the uploaded files are listed below.}

\subsection{Notebook-to-Script Conversion}
Real-world notebooks are full of Jupyter-only constructs (\texttt{!pip
install}, cell magics, shell commands) that crash outside Jupyter. At upload
time, the platform converts every notebook into a clean, runnable script
(Figure~\ref{fig:notebook-convert}): cells that already parse as pure Python
are kept byte-for-byte; the others are translated with IPython's own
transformer~\cite{nbconvert_docs}, then rewritten into portable Python ---
\texttt{!pip install x} becomes a real \texttt{subprocess} install (and the
packages are reported to the interface), \texttt{\%\%time} keeps its body,
display-only magics are commented out. The final script must compile, or the
system falls back safely to in-pod conversion.

\begin{figure}[H]\centering
  \includegraphics[width=\textwidth,keepaspectratio]{notebook_convert}
  \caption{Notebook-to-script conversion: before and after}
  \label{fig:notebook-convert}
\end{figure}

\subsection{The Training Run}
A training run executes inside an isolated pod, following the training path
designed in Section~\ref{sec:pipeline-org}: download code and dataset, install
missing dependencies (with version pins for known conflicts), and run the
entry script under a generated runner that enables automatic tracking first.
The runner also neutralizes blocking calls (\texttt{plt.show()},
\texttt{input()}) and shims runtime-only modules such as \texttt{google.colab},
so notebooks exported from cloud environments run unchanged. While the run
executes, the backend streams the pod's logs to the browser in real time
(Figure~\ref{fig:screen_logs}).

\bigshot{screen_logs}{The live log terminal}{The training pod's output streams
to the browser as the run executes, so the user follows dependency
installation, notebook conversion, and training in real time.}

\subsection{Experiments Intelligence}
Captured runs are made comparable: the experiments view groups runs by model
family, ranks them on \emph{evaluation} metrics only (never training scores),
flags the champion model, and registers any run as a model version in one
click (Figures~\ref{fig:expintel} and~\ref{fig:screen_experiments}).

\begin{figure}[H]\centering
  \includegraphics[width=\textwidth,keepaspectratio]{expintel_features}
  \caption{Experiments intelligence: grouping, fair ranking, champion, one-click register}
  \label{fig:expintel}
\end{figure}

\bigshot{screen_experiments}{The experiments view}{Runs are grouped by model
family with their captured metrics, the champion run is highlighted, and any
run can be registered as a model version in one click.}

\subsection{Run Details and Artifact Browser}
\label{sec:artifact-browser}
Each run opens into a detail view: a status header, headline \textbf{metric
cards}, and a searchable parameter list. Its second half is an
\textbf{artifact browser}: the backend lists the run's files directly through
the \gls{s3} \gls{api} and streams any file with its correct content type, so
evaluation plots render inline, folders are navigable, and every file is
downloadable. Both endpoints resolve the run back to a project owned by the
caller and answer \texttt{404} otherwise, so artifacts are never leaked across
users.

\subsection{Sprint Review}
The review demonstrated the full training loop on real notebooks taken from
Colab, not simplified examples. The retrospective identified dependency
conflicts inside training pods as the most frequent failure cause, which
motivated the automatic missing-dependency installation and, later, the
pre-flight code analyzer of Sprint~4.

% ===========================================================================
\section{Sprint 3 --- Registry and Deployment}
\label{sec:sprint3}
The third sprint turns tracked runs into deployable assets: a registry with
lifecycle stages, one-click deployment, and an in-interface prediction tester
(Figure~\ref{fig:sprint3-strip}).

\begin{figure}[H]\centering
  \includegraphics[width=\textwidth,keepaspectratio]{sprint3_features}
  \caption{Sprint 3 at a glance: registry, promotion, KServe deployment, tester}
  \label{fig:sprint3-strip}
\end{figure}

\begin{table}[H]\centering
\renewcommand{\arraystretch}{1.15}
\caption{Sprint~3 backlog}\label{tab:backlog-s3}
\small
\begin{tabularx}{\textwidth}{|c|X|c|}
\hline
\rowcolor{primary!15}
\textbf{\#} & \textbf{Task} & \textbf{Pts}\\
\hline
8  & Register model versions from successful runs & 2\\ \hline
8  & Build the model registry interface & 3\\ \hline
9  & Promote models across Staging/Production/Archived & 3\\ \hline
10 & Deploy a model as an InferenceService on the cluster & 5\\ \hline
10 & Detect the model's framework and serve XGBoost natively & 3\\ \hline
10 & Track deployment status and endpoint & 3\\ \hline
11 & Build the in-UI prediction tester & 5\\ \hline
11 & Surface readable model-server errors to the caller & 2\\ \hline
\end{tabularx}
\end{table}

\subsection{Registry and Promotion}
Successful runs can be registered as named, versioned models. Each version
carries a lifecycle stage (Staging, Production, Archived) that the user
changes from the interface; the tracking server remains the source of truth,
mirrored in the platform's database (Figure~\ref{fig:screen_registry}).

\bigshot{screen_registry}{The model registry}{Each registered version shows its
lifecycle-stage badge (Staging / Production / Archived) and can be promoted or
deployed directly from this view.}

\subsection{One-Click Deployment}
Deploying a model reduces to creating one declarative \texttt{InferenceService}
resource, following the serving path of Section~\ref{sec:pipeline-org}; the
operator reconciles it into a running model server. The backend then polls the
resource's status conditions and updates the deployment record until the
service is \texttt{READY}, so the interface always reflects the real cluster
state (Figure~\ref{fig:screen_deploy}).

\bigshot{screen_deploy}{The deployment panel}{The deployment progresses through
its status conditions to \texttt{READY}, and the live prediction endpoint of
the served model is displayed.}

\subsection{Multi-Framework Serving}
\label{sec:multi-framework}
The serving operator provides one model server per framework, and serving an
XGBoost model through the scikit-learn server fails at startup. The backend
therefore \textbf{detects the framework} at deployment time by listing the
file names in the model's artifact folder: a booster file
(\texttt{model.bst} / \texttt{.json} / \texttt{.ubj} / \texttt{.xgb}) marks an
XGBoost model, anything else defaults to scikit-learn. One incompatibility
remains: MLflow saves XGBoost boosters as \texttt{model.xgb}, a name the
XGBoost server refuses. In that case the backend \emph{materializes} a servable
copy (\texttt{model.xgb} $\rightarrow$ \texttt{model.bst}) in the object store
before creating the service. The user never sees any of this: they click
``Deploy'' and the right runtime serves the right file.

\subsection{Prediction and Error Handling}
\label{sec:prediction-errors}
The backend forwards prediction requests to the model server through the
cluster's \gls{api}, so no model needs a public address; the interface offers
a tester with sample inputs (Figure~\ref{fig:screen_predict}). Failures are
normalized into a dedicated \texttt{PredictionError} with a \emph{readable}
message: when the model server rejects an input, its own error text is returned
with \texttt{502}; when it does not answer in time, the request fails fast with
\texttt{504}. Forwarded predictions run in a worker thread, so a slow model
server never blocks the asynchronous event loop.

\bigshot{screen_predict}{The prediction tester}{The user submits sample inputs
to the deployed model and the returned prediction is displayed, with a readable
error message when an input is malformed.}

\subsection{Sprint Review}
The review showed a model going from a tracked run to a live, tested service
in a few clicks. The retrospective surfaced the framework mismatch between
MLflow's artifacts and the serving runtimes; the automatic framework detection
and booster materialization presented above were the direct outcome.

% ===========================================================================
\section{Sprint 4 --- Observability, Public API, and Administration}
\label{sec:sprint4}
The final sprint makes the platform observable, opens deployed models to the
outside world, and adds administration and diagnostics
(Figure~\ref{fig:sprint4-strip}).

\begin{figure}[H]\centering
  \includegraphics[width=\textwidth,keepaspectratio]{sprint4_features}
  \caption{Sprint 4 at a glance: dashboards, telemetry, API keys, administration}
  \label{fig:sprint4-strip}
\end{figure}

\begin{table}[H]\centering
\renewcommand{\arraystretch}{1.15}
\caption{Sprint~4 backlog}\label{tab:backlog-s4}
\small
\begin{tabularx}{\textwidth}{|c|X|c|}
\hline
\rowcolor{primary!15}
\textbf{\#} & \textbf{Task} & \textbf{Pts}\\
\hline
12 & Collect cluster and per-pod resource metrics & 3\\ \hline
12 & Build the monitoring dashboard & 2\\ \hline
12 & Record per-prediction telemetry and aggregate serving KPIs & 3\\ \hline
13 & Generate and store per-deployment API keys (hashed) & 3\\ \hline
13 & Expose a public, key-authenticated prediction endpoint & 2\\ \hline
14 & Capture and surface user-script errors for failed runs & 3\\ \hline
14 & Add a static pre-flight code analyzer with warnings & 2\\ \hline
15 & Add the administrator role and administration area & 5\\ \hline
\end{tabularx}
\end{table}

\subsection{Cluster Observability}
The backend queries the cluster's metrics service~\cite{metrics_server} for
per-pod \gls{cpu} and memory usage and aggregates cluster-wide figures,
surfaced as live gauges and charts on the dashboard
(Figure~\ref{fig:screen_dashboard}).

\bigshot{screen_dashboard}{The metrics dashboard}{Cluster-wide and per-pod
\gls{cpu} and memory usage are shown as live gauges and charts, so the health
of the cluster and of each deployed model is visible at a glance.}

\subsection{Serving Telemetry and KPIs}
\label{sec:serving-telemetry}
Resource metrics say whether a model server is \emph{running}, not how it is
\emph{behaving}. The platform therefore records one \texttt{PredictionLog} row
for every served prediction --- from the tester or the public \gls{api} ---
with latency, \gls{http} outcome, batch size, and origin. Rows are written
best-effort and contain no payloads.

A per-project endpoint aggregates these rows over a selectable window into the
\glspl{kpi} that matter once a model is live: request volume, error rate,
average and 95\textsuperscript{th}-percentile latency, and a bucketed timeline
rendered as charts in the Monitoring tab.

\subsection{Public Prediction API with API Keys}
Deployment makes a model reachable from inside the platform; this sprint opens
it to the outside world. The owner creates \textbf{per-deployment \gls{api}
keys}: the key (\texttt{mlops\_<random>}) is shown \textbf{once}, only its
\gls{sha}-256 hash and a display prefix are stored, and a key can be revoked
at any time.

External clients call a dedicated public endpoint, outside the
\gls{jwt}-protected \gls{api}, with the key in the \texttt{Authorization}
header; the backend hashes and looks it up, confirms the deployment is ready,
and forwards the request (Figure~\ref{fig:public-api}). Every public call also
lands in the serving telemetry, so external usage is visible in the same
Monitoring tab. A snippet generator produces ready-to-paste cURL, JavaScript,
and Python examples (Figures~\ref{fig:public-api} and~\ref{fig:screen_apikeys}).

\begin{figure}[H]\centering
  \includegraphics[width=0.9\textwidth,height=0.42\textheight,keepaspectratio]{public_api}
  \caption{Public API flow: key issuance, authenticated request, response}
  \label{fig:public-api}
\end{figure}

\bigshot{screen_apikeys}{The API-keys panel}{Per-deployment keys are created and
revoked here (only their hash is stored), and the code-snippet generator
produces ready-to-paste cURL, JavaScript, and Python calls.}

\subsection{Administration}
A \texttt{role} field distinguishes users from administrators; a backend
dependency guards every administrative endpoint and an Angular guard protects
the routes. The administration area provides user management (activate,
deactivate, promote, demote, purge, delete, with anti-lockout safeguards), a
platform overview, and read-only oversight of all users' resources. Because the
role is read from the database on every request, a demotion takes effect
immediately (Figure~\ref{fig:screen_admin}).

\bigshot{screen_admin}{The administration area}{Every account is listed with its
role and status; an administrator can activate, deactivate, promote, demote,
purge, or delete a user, and oversee all users' resources.}

\subsection{Diagnostics and Robustness}
The platform captures the user-script error of a failed run and stores it
durably, so the interface shows a concise ``why did this fail'' explanation. A
static pre-flight analyzer also inspects uploaded code before launch and warns
about common pitfalls: missing training calls, hard-coded local paths,
notebook-only constructs (Figure~\ref{fig:screen_diagnostics}).

\bigshot{screen_diagnostics}{The failure-explanation panel}{The captured
user-script error is surfaced on the failed run, telling the user \emph{why} it
failed without digging through raw job logs.}

\subsection{Sprint Review}
The final review demonstrated an external client calling a deployed model with
an \gls{api} key and the resulting traffic appearing live in the Monitoring
tab. With administration and diagnostics accepted, the product backlog of
Chapter~\ref{chap:context} was fully delivered.

% ===========================================================================
\section{Continuous Deployment and Environments}
\label{sec:cicd}
The platform itself is delivered like software. Every push to the main branch
triggers a GitHub Actions workflow that builds the backend and frontend
images with registry-backed layer caching, pushes them to \gls{acr}, and
upgrades the Helm release on the \gls{aks} cluster. A concurrency group cancels
any in-progress run when a newer commit arrives, so the cluster always
converges on the latest version of the main branch.

The same Helm chart deploys the full stack to both environments, with only a
values file changing (Table~\ref{tab:environments}).

\begin{table}[H]\centering
\renewcommand{\arraystretch}{1.2}
\caption{Execution environments}\label{tab:environments}
\small
\begin{tabularx}{\textwidth}{|>{\raggedright\arraybackslash}p{3.4cm}|X|X|}
\hline
\rowcolor{primary!15}
\textbf{Aspect} & \textbf{Development (KinD)} & \textbf{Production (AKS)}\\
\hline
Cluster & Kubernetes-in-Docker, local & Azure Kubernetes Service, managed\\ \hline
Purpose & Development and tests & Live, public platform\\ \hline
Images & Loaded locally & Pulled from \gls{acr} by CI/CD\\ \hline
Public access & None (localhost) & Cloudflare Tunnel + managed \gls{tls}\\ \hline
Manifests & \multicolumn{2}{c|}{Same Helm chart, different values file}\\ \hline
\end{tabularx}
\end{table}

% ---------------------------------------------------------------------------
\section{Conclusion}
Four sprints turned the design into a complete platform: identity and
projects, a training-and-tracking engine with notebook conversion, a registry
with one-click multi-framework deployment, and observability with a public
\gls{api} and administration --- all continuously delivered to a production
cluster. The next chapter validates the result, scenario by scenario and on a
real use case.
